\documentclass[12pt, a4paper]{article}

% Packages
\usepackage[margin=1in]{geometry}
\usepackage{booktabs}
\usepackage{hyperref}
\usepackage{enumitem}
\usepackage{setspace}
\usepackage[numbers]{natbib}
\usepackage{array}
\usepackage{tabularx}

% Formatting
\onehalfspacing

\hypersetup{
    colorlinks=true,
    linkcolor=blue,
    citecolor=blue,
    urlcolor=blue
}

\begin{document}

% ============ HEADER ============
\begin{center}
    {\large \textbf{National Institute of Technology Calicut}}\\[0.2cm]
    {\normalsize Department of Computer Science \& Engineering}\\[0.8cm]
    
    {\Large \textbf{Proposed Research Topic}}\\[0.4cm]
    
    {\large \textbf{Multimodal Anomaly Detection in Industrial Quality Inspection\\ using Vision-Language Models}}\\[0.8cm]
    
    \begin{tabular}{rl}
        \textbf{Student:} & [Your Name] \\
        \textbf{Roll No:} & [Your Roll Number] \\
        \textbf{Guide:} & Prof.\ [Professor Name] \\
        \textbf{Date:} & March 2, 2026 \\
    \end{tabular}
\end{center}

\vspace{0.5cm}
\hrule
\vspace{0.5cm}

% ============ 1. AREA OF INTEREST ============
\section{Area of Interest}

\textbf{Multimodal anomaly detection} involves using multiple types of data (e.g., images and text) together to identify items that deviate from expected normal patterns. In industrial manufacturing, this translates to automated quality inspection --- detecting product defects such as scratches, dents, missing components, or structural damage on production lines.

Current anomaly detection methods predominantly rely on \textbf{single-modality} (image-only) approaches and require training on normal samples for every new product category, which limits their scalability. Recent advances in \textbf{Vision-Language Models} (VLMs) offer a promising alternative by jointly understanding images and natural language, potentially enabling \textbf{zero-shot} anomaly detection --- detecting defects without any product-specific training.

% ============ 2. INITIAL FINDINGS ============
\section{Initial Findings from Literature}

I studied several recent works in this area. Here are the key findings from five representative papers spanning traditional baselines, CLIP-based methods, MLLM-based approaches, and benchmarking:

\subsection{Paper 1: WinCLIP \cite{jeong2023winclip}}
\begin{itemize}[leftmargin=*]
    \item \textbf{What they did:} Applied the CLIP vision-language model for zero-shot anomaly detection in manufacturing. Used text prompts like ``a photo of a damaged [product]'' vs ``a photo of a normal [product]'' and a sliding window approach for defect localization.
    \item \textbf{Key result:} Achieved 91.8\% image-level AUROC on the MVTec AD benchmark \textbf{without any training}, competitive with many supervised methods.
    \item \textbf{Limitation:} Relies on hand-crafted prompts; cannot explain \textit{why} something is anomalous.
    \item \textbf{Code:} \url{https://github.com/caoyunkang/WinClip}
\end{itemize}

\subsection{Paper 2: AnomalyGPT \cite{gu2024anomalygpt}}
\begin{itemize}[leftmargin=*]
    \item \textbf{What they did:} Used a Large Vision-Language Model (LVLM) fine-tuned on simulated anomaly data for industrial defect detection. The model can detect anomalies, localize them, and describe them in natural language through dialogue.
    \item \textbf{Key result:} Achieves state-of-the-art results on MVTec AD with the added ability to \textit{explain} detected defects in words.
    \item \textbf{Limitation:} Requires fine-tuning on simulated data; not truly zero-shot.
    \item \textbf{Code:} \url{https://github.com/CASIA-IVA-Lab/AnomalyGPT}
\end{itemize}

\subsection{Paper 3: MMAD Benchmark \cite{jiang2025mmad}}
\begin{itemize}[leftmargin=*]
    \item \textbf{What they did:} Created a comprehensive benchmark to evaluate Multimodal LLMs (GPT-4o, Claude, open-source models) on industrial anomaly detection across 7 subtasks.
    \item \textbf{Key result:} Even the best models (GPT-4o) show \textbf{significant room for improvement}, indicating this is an open and active research problem.
    \item \textbf{Takeaway:} There is a clear opportunity to develop better MLLM-based anomaly detection approaches.
    \item \textbf{Code:} \url{https://github.com/jam-cc/MMAD}
\end{itemize}

\subsection{Paper 4: AnomalyCLIP \cite{zhou2024anomalyclip}}
\begin{itemize}[leftmargin=*]
    \item \textbf{What they did:} Introduced object-agnostic learnable text prompts for CLIP-based zero-shot anomaly detection. Instead of hand-crafting prompts per product category, AnomalyCLIP optimizes continuous prompt embeddings that capture universal ``normal'' vs ``anomalous'' semantics across all object types, while keeping the CLIP backbone frozen.
    \item \textbf{Key result:} Achieves approximately 94\% image-level AUROC on MVTec AD and demonstrates significant improvement over WinCLIP on cross-domain evaluation with unseen product categories.
    \item \textbf{Limitation:} Still requires a training phase for prompt optimization; the learned prompts may overfit to the training categories.
    \item \textbf{Code:} \url{https://github.com/zqhang/AnomalyCLIP}
\end{itemize}

\subsection{Paper 5: PatchCore \cite{roth2022patchcore}}
\begin{itemize}[leftmargin=*]
    \item \textbf{What they did:} Proposed a memory-bank-based approach for industrial anomaly detection. Extracts patch-level features from a pre-trained CNN, stores representative normal features via coreset subsampling, and detects anomalies using nearest-neighbor distance --- test patches far from all stored normal patches are flagged as anomalous.
    \item \textbf{Key result:} Achieves \textbf{99.1\% image-level AUROC} on MVTec AD --- the strongest traditional baseline and current upper-bound for non-VLM methods.
    \item \textbf{Limitation:} Requires category-specific normal training images; cannot generalize to unseen products and provides no natural language explanations of defects.
    \item \textbf{Code:} \url{https://github.com/amazon-science/patchcore-inspection}
\end{itemize}

% ============ 3. IDENTIFIED GAP ============
\section{Preliminary Research Gap}

From these five papers, I observe:
\begin{enumerate}[leftmargin=*]
    \item CLIP-based methods (WinCLIP) enable zero-shot scoring but rely on hand-crafted prompts and cannot explain anomalies.
    \item Learnable prompt approaches (AnomalyCLIP) improve generalization but still require a training phase for prompt optimization.
    \item MLLM-based methods (AnomalyGPT) can explain anomalies but require fine-tuning on simulated data.
    \item Current MLLMs (as shown by MMAD benchmark) still perform poorly out-of-the-box for industrial anomaly detection.
    \item Traditional methods (PatchCore, 99.1\%) remain the upper-bound, but require category-specific training data and offer no interpretability.
\end{enumerate}

\textbf{Gap:} A systematic study of how to effectively use Vision-Language Models and open-source MLLMs (e.g., LLaVA) together for zero-shot anomaly detection --- combining quantitative scoring with natural language reasoning --- is currently lacking. Furthermore, no work has systematically studied the impact of prompt engineering strategies on CLIP-based anomaly detection performance.

% ============ 4. PROPOSED DIRECTION ============
\section{Proposed Research Direction}

I would like to investigate:

\begin{quote}
\textit{``How can pre-trained Vision-Language Models and Multimodal Large Language Models be effectively leveraged for zero-shot anomaly detection in industrial quality inspection, achieving robust defect detection and interpretable explanations without task-specific training?''}
\end{quote}

\textbf{Proposed approach:}
\begin{itemize}[leftmargin=*]
    \item \textbf{Stage 1 --- VLM-based anomaly scoring:} Use CLIP with structured prompt engineering to compute anomaly scores at image and pixel level. Explore prompt variations and compositional ensembles for robust detection. Systematically compare hand-crafted prompts (WinCLIP-style) vs structured prompt strategies to understand the impact of prompt design.
    \item \textbf{Stage 2 --- MLLM-based reasoning:} For images flagged as anomalous, use an open-source MLLM (LLaVA) to generate natural language descriptions of the detected defect --- its type, location, and severity.
    \item \textbf{Stage 3 --- Systematic evaluation:} Benchmark this two-stage hybrid pipeline against traditional baselines (PatchCore, PaDiM) and existing VLM methods (WinCLIP, AnomalyCLIP) on MVTec AD and VisA datasets using the evaluation framework inspired by MMAD.
\end{itemize}

The novelty lies in the \textbf{hybrid two-stage pipeline} (CLIP for scoring + MLLM for reasoning) and the \textbf{systematic prompt engineering study} --- neither of which has been explored in prior work.

% ============ 5. PAPERS FOR FURTHER REVIEW ============
\section{Papers for Further Review}

In addition to the five papers studied above, I plan to review the following works along with their publicly available code for a deeper understanding and implementation reference:

\begin{table}[h]
\centering
\small
\renewcommand{\arraystretch}{1.4}
\begin{tabularx}{\textwidth}{@{}p{0.7cm}p{5.5cm}p{1.2cm}X@{}}
\toprule
\textbf{No.} & \textbf{Paper} & \textbf{Venue} & \textbf{Code Repository} \\
\midrule
1 & LLaVA: Visual Instruction Tuning \cite{liu2023llava} & NeurIPS 2023 & \url{https://github.com/haotian-liu/LLaVA} \\
\addlinespace
2 & Anomalib: A Deep Learning Library for Anomaly Detection \cite{anomalib2022} & -- & \url{https://github.com/openvinotoolkit/anomalib} \\
\addlinespace
3 & VCP-CLIP: Visual Context Prompting for Zero-Shot Anomaly Segmentation & ECCV 2024 & \url{https://github.com/xiaozhi-qu/VCP-CLIP} \\
\addlinespace
4 & VELM: Vision Expert-Enhanced LLM for Anomaly Detection & 2024 & -- \\
\bottomrule
\end{tabularx}
\caption{Additional papers and code repositories for review}
\end{table}

\textbf{Purpose of reviewing these works:}
\begin{itemize}[leftmargin=*]
    \item \textbf{LLaVA} --- The open-source MLLM I plan to use for Stage 2 (defect reasoning and explanation). I will study its architecture and inference pipeline.
    \item \textbf{Anomalib} --- Intel's anomaly detection library implementing multiple baselines (PatchCore, PaDiM, etc.). This will serve as the experimental framework for running and comparing traditional methods efficiently.
    \item \textbf{VCP-CLIP} --- Recent ECCV 2024 work on visual context prompting for anomaly segmentation, relevant to understanding visual prompting strategies.
    \item \textbf{VELM} --- Hybrid approach combining traditional anomaly detection methods (as vision experts) with LLMs, closely related to my two-stage pipeline concept.
\end{itemize}

% ============ 6. NEXT STEPS ============
\section{Next Steps}

\begin{itemize}[leftmargin=*]
    \item Review the above papers and their codebases in detail
    \item Set up the experimental environment (Python, PyTorch, Anomalib, OpenCLIP)
    \item Run preliminary experiments with CLIP on MVTec AD to establish initial baselines
    \item Formalize the problem statement and detailed methodology
    \item Prepare a comprehensive literature survey for mid-semester review
\end{itemize}

% ============ REFERENCES ============
\begin{thebibliography}{10}

\bibitem{jeong2023winclip}
J.~Jeong, Y.~Zou, T.~Kim, D.~Zhang, A.~Ravichandran, and O.~Dabeer.
\newblock WinCLIP: Zero-/few-shot anomaly classification and segmentation.
\newblock In \emph{CVPR}, 2023.

\bibitem{gu2024anomalygpt}
Z.~Gu et~al.
\newblock AnomalyGPT: Detecting industrial anomalies using large vision-language models.
\newblock In \emph{AAAI}, 2024.

\bibitem{jiang2025mmad}
X.~Jiang et~al.
\newblock MMAD: Multi-modal anomaly detection benchmark.
\newblock In \emph{ICLR}, 2025.

\bibitem{zhou2024anomalyclip}
Q.~Zhou et~al.
\newblock AnomalyCLIP: Object-agnostic prompt learning for zero-shot anomaly detection.
\newblock In \emph{ICLR}, 2024.

\bibitem{roth2022patchcore}
K.~Roth, L.~Pemula, J.~Zepeda, B.~Sch{\"o}lkopf, T.~Brox, and P.~Gehler.
\newblock Towards total recall in industrial anomaly detection.
\newblock In \emph{CVPR}, 2022.

\bibitem{liu2023llava}
H.~Liu, C.~Li, Q.~Wu, and Y.~J.~Lee.
\newblock Visual instruction tuning.
\newblock In \emph{NeurIPS}, 2023.

\bibitem{anomalib2022}
S.~Akcay et~al.
\newblock Anomalib: A deep learning library for anomaly detection.
\newblock In \emph{IEEE ICIP}, 2022.

\end{thebibliography}

\end{document}
